\appendix
\section{The twenty-four Kummer expressions}\label{app:kummer24}

This appendix records a complete working list of the twenty-four classical
expressions generated from one local hypergeometric solution by the six
fractional-linear arguments
\[
 z,\quad 1-z,\quad z^{-1},\quad (1-z)^{-1},\quad
 \frac{z}{z-1},\quad 1-z^{-1}
\]
and the elementary powers of \(z\), \(1-z\), and \(-z\).  They are grouped
into six quartets.  Within a simply connected rational polygonal domain on
which all required powers have certified logarithm branches, each expression
is a solution of the same pulled-back differential equation.  Equality of two
expressions is therefore reduced, as in \cref{thm:formal-uniqueness}, to equality
of two initial coefficients.  This is a \StatusB{} theorem schema.  Without
branch certificates, the table is only a formal list of local expressions;
no global single-valued equality is intended.

For readability write \(F(a,b;c;w)=\F{a}{b}{c}{w}\).  A convenient set of
representatives is
\begin{align*}
K_{1,1}&=F(a,b;c;z),\\
K_{1,2}&=(1-z)^{c-a-b}F(c-a,c-b;c;z),\\
K_{1,3}&=(1-z)^{-a}F\left(a,c-b;c;\frac{z}{z-1}\right),\\
K_{1,4}&=(1-z)^{-b}F\left(c-a,b;c;\frac{z}{z-1}\right);
\end{align*}
\begin{align*}
K_{2,1}&=F(a,b;a+b+1-c;1-z),\\
K_{2,2}&=z^{1-c}F(a+1-c,b+1-c;a+b+1-c;1-z),\\
K_{2,3}&=z^{-a}F(a,a+1-c;a+b+1-c;1-z^{-1}),\\
K_{2,4}&=z^{-b}F(b+1-c,b;a+b+1-c;1-z^{-1});
\end{align*}
\begin{align*}
K_{3,1}&=(-z)^{-a}F(a,a+1-c;a+1-b;z^{-1}),\\
K_{3,2}&=(-z)^{b-c}(1-z)^{c-a-b}
 F(1-b,c-b;a+1-b;z^{-1}),\\
K_{3,3}&=(1-z)^{-a}F\left(a,c-b;a+1-b;(1-z)^{-1}\right),\\
K_{3,4}&=(-z)^{1-c}(1-z)^{c-a-1}
 F\left(a+1-c,1-b;a+1-b;(1-z)^{-1}\right);
\end{align*}
\begin{align*}
K_{4,1}&=(-z)^{-b}F(b+1-c,b;b+1-a;z^{-1}),\\
K_{4,2}&=(-z)^{a-c}(1-z)^{c-a-b}
 F(1-a,c-a;b+1-a;z^{-1}),\\
K_{4,3}&=(1-z)^{-b}F\left(b,c-a;b+1-a;(1-z)^{-1}\right),\\
K_{4,4}&=(-z)^{1-c}(1-z)^{c-b-1}
 F\left(b+1-c,1-a;b+1-a;(1-z)^{-1}\right);
\end{align*}
\begin{align*}
K_{5,1}&=z^{1-c}F(a+1-c,b+1-c;2-c;z),\\
K_{5,2}&=z^{1-c}(1-z)^{c-a-b}F(1-a,1-b;2-c;z),\\
K_{5,3}&=z^{1-c}(1-z)^{c-a-1}
 F\left(a+1-c,1-b;2-c;\frac{z}{z-1}\right),\\
K_{5,4}&=z^{1-c}(1-z)^{c-b-1}
 F\left(b+1-c,1-a;2-c;\frac{z}{z-1}\right);
\end{align*}
\begin{align*}
K_{6,1}&=(1-z)^{c-a-b}F(c-a,c-b;c+1-a-b;1-z),\\
K_{6,2}&=z^{1-c}(1-z)^{c-a-b}
 F(1-a,1-b;c+1-a-b;1-z),\\
K_{6,3}&=z^{a-c}(1-z)^{c-a-b}
 F(c-a,1-a;c+1-a-b;1-z^{-1}),\\
K_{6,4}&=z^{b-c}(1-z)^{c-a-b}
 F(c-b,1-b;c+1-a-b;1-z^{-1}).
\end{align*}

The first quartet consists of the identity, Euler's transformation, and the
two Pfaff transformations.  The remaining quartets are obtained by changing
the distinguished singular point and choosing either local exponent there.
At resonant parameters some displayed pairs coalesce or their normalizations
become singular.  The correct replacement is then a logarithmic Frobenius
solution, so the generic twenty-four-label statement must not be read as
asserting twenty-four linearly independent functions.

\section{Primitive recursive precision schedules}\label{app:schedules}

The arguments in the main text repeatedly use the same small collection of
precision compilers.  This appendix makes them explicit enough to serve as an
implementation specification.

\subsection{Geometric domination}

Let \(0\le r<1\) be rational and let \(A\ge0\).  Put
\[
 d=\operatorname{den}(1-r),\qquad
 E(q,A,r)=\min\left\{N\le B:
 A r^N\le 2^{-q}(1-r)\right\},
\]
where a primitive recursive search bound is
\[
 B=(q+\sigma(A)+2)d,
 \qquad
 \sigma(A)=\min\{s\le \lceil A\rceil+1:2^s\ge\max(1,A)\}.
\]
The elementary inequality \((1-1/d)^d\le 1/2\), proved after clearing
denominators by induction on \(d\), shows that the bounded search succeeds
after enlarging \(B\) by at most a fixed linear factor.  Consequently
\[
 \sum_{n\ge E(q,A,r)} A r^n\le 2^{-q}.
\]
All occurrences of a phrase such as ``choose \(N\) so that the geometric tail
is below \(2^{-q}\)'' can be replaced by this bounded search.

\subsection{Hypergeometric ratio schedule}

For fixed rational-complex parameters with denominator margins
\(\norm{c+n}\ge\eta(n)>0\), the exact ratio is
\[
 R_n(z)=\frac{(a+n)(b+n)}{(c+n)(n+1)}z.
\]
After clearing the parameter denominators, the inequality
\(\norm{R_n(z)}\le r<1\) is a polynomial inequality in \(n\).  The proof of
\cref{lem:ratio-cutoff} supplies an explicit polynomial bound \(N_0\) beyond
which it holds.  The complete cutoff is
\[
 N(q)=\max\bigl(N_0,E(q,A,r)\bigr),
\]
where \(A\) bounds the normalized \(N_0\)-th term.  This is primitive
recursive in the input codes and requested precision.

\subsection{Endpoint truncation for Euler's integral}

Assume rational witnesses
\[
 0<\delta_b\le\RePart b,
 \qquad
 0<\delta_{c-b}\le\RePart(c-b).
\]
For an integrand bounded by
\(M t^{\delta_b-1}(1-t)^{\delta_{c-b}-1}\), choose dyadic
\(\varepsilon=2^{-m}\) by bounded search so that
\[
 \frac{M}{\delta_b}\varepsilon^{\delta_b}
 +\frac{M}{\delta_{c-b}}\varepsilon^{\delta_{c-b}}
 \le 2^{-q-2}.
\]
Rational powers are evaluated by the certified exponential/logarithm names of
\cref{sec:arithmetic}.  On \([\varepsilon,1-\varepsilon]\), use the supplied
modulus of continuity to choose a mesh whose Riemann error is at most
\(2^{-q-2}\), and evaluate each node to precision \(2^{-q-2}/(M+1)\).  The
three errors add to at most \(2^{-q}\).  Every search is bounded by a function
of the displayed rational data.

\subsection{Path continuation schedule}

For a polygonal path \(\gamma\) avoiding \(\{0,1\}\), a continuation
certificate contains rational margins
\[
 \rho_j\le \min_{z\in S_j}\bigl(\abs z,\abs{1-z}\bigr).
\]
After decomposing the path into segments \(S_j\).  On each segment the ODE is
rewritten as a first-order system with a rational Lipschitz bound \(L_j\).
A truncated Taylor or Picard step of length \(h_j\) is accepted when
\(L_jh_j\le 1/4\).  The order is chosen by the geometric compiler above,
and the local errors are budgeted as
\(2^{-q-j-2}\).  Finite summation of the local budgets proves the final
\(2^{-q}\) error.  This realizes continuation without a compactness or
choice principle.

\section{Object-theory target schemas}\label{app:targets}

A theorem in the present report becomes a closed first-order iPRA sentence
only after the primitive recursive programs naming its parameters and
certificates have been fixed.  The following patterns are the intended
translation targets in ProveIt's \code{Prf PRAAx} calculus.

\subsection{Exact coefficient identity}

If \(A(n)\) and \(B(n)\) are integer-coded rational expressions, equality is
represented after cross multiplication by a natural-number equation
\[
 \operatorname{num}A(n)\operatorname{den}B(n)
 =\operatorname{num}B(n)\operatorname{den}A(n).
\]
The target is the universal closure of this quantifier-free equation.  A
simultaneous recurrence can be encoded by the pairing machinery already
available in the repository.

\subsection{Equality of effective complex numbers}

For primitive recursive complex-rational names \(X(q)\) and \(Y(q)\), the
statement \(X=Y\) is compiled to the schema
\[
 \forall q\;
 2^q\,\norm{X(q+2)-Y(q+2)}\le 1.
\]
The norm is the rational \(\ell^1\)-norm, so the displayed predicate is a
finite integer inequality.  The two lost bits absorb the independent naming
errors.  Uniform identities add the bounded parameter code as another free
number variable.

\subsection{Convergence with a modulus}

For partial sums \(S(N)\) and a primitive recursive modulus \(M(q)\), Cauchy
convergence is the arithmetical formula
\[
 \forall q\,\forall m,n\;
 \bigl(M(q)\le m\wedge M(q)\le n\bigr)
 \to 2^q\norm{S(m)-S(n)}\le1.
\]
When the proof uses a geometric tail, the implication is discharged by the
exact finite geometric identity and the cutoff of \cref{app:schedules}.

\subsection{Apartness and branches}

A guarded division \(x/y\) carries an integer witness \(k\) satisfying
\(2^{-k}\le\abs y\).  A branch certificate for \(z^\alpha\) is a finite list
of rational polygons, overlap points, logarithm names, and estimates asserting
that the image stays a rational positive distance from the selected cut.  All
checks are bounded conjunctions of integer inequalities.  Thus the apparent
higher-type data are replaced by one number coding a finite certificate and a
primitive recursive checker.

\subsection{Suggested closed milestone theorems}

The first object-theory implementation should target, in order:
\begin{enumerate}
\item the Pochhammer recurrence and zero criterion;
\item the coefficient ratio, differentiation identity, and ODE residual;
\item terminating Chu--Vandermonde and Jacobi/Legendre specializations;
\item the three basic contiguous relations and their finite continued-fraction
      eliminant;
\item the cubic radical coefficient recurrence;
\item the generic disk tail theorem for a fixed rational parameter tuple;
\item Euler's integral for one fixed positive rational parameter tuple;
\item Gauss summation in the same tuple, followed by a parameterized schema.
\end{enumerate}
The first five use only finite algebra.  The last three exercise the effective
analysis interfaces without yet requiring branch continuation.

\section{Archive map and reproducibility}\label{app:archive}

\begin{center}
\begingroup\small
\begin{tabularx}{\textwidth}{@{}p{0.35\textwidth}X@{}}
\toprule
Path & Role\\
\midrule
\path{report.tex}, \path{report.pdf} & Complete paper and compiled rendering.\\
\path{sections/} & Source sections, including the generated inventory table.\\
\path{theorem_inventory.csv} & Source-row coverage ledger with status, route, and qualifications.\\
\path{lean/HypergeometricFinite.lean} & Executable-looking finite algebra layer and theorem statements; uncompiled outline.\\
\path{lean/HypergeometricIPRA.lean} & Object-theory translation skeleton aimed at \code{Prf PRAAx}.\\
\path{lean/FORMALIZATION_PLAN.md} & Dependency graph, induction invariants, and completion checklist.\\
\path{code/validate_finite.py} & Exact rational and high-precision numerical regression checks.\\
\path{evidence/validation.json} & Deterministic output of the validation program.\\
\path{SOURCE_MANIFEST.md} & Web and repository sources, scope, and inspected revision.\\
\path{SHA256SUMS.txt} & Checksums for all packaged files except the checksum file itself.\\
\bottomrule
\end{tabularx}
\endgroup
\end{center}

The PDF is rebuilt with \code{make pdf}; the finite evidence with
\code{make validate}; and both plus checksums with \code{make all}.  The
validation program proves no general theorem, but it is useful for detecting
sign, index-shift, and normalization mistakes in the human and Lean plans.
